\section{制度不是命令的别名}
\label{sec:synthesis-institutions}

两晋南北朝的制度史常被压成几串名词：西晋律令、九品中正、侨置与土断、北魏均田三长、府兵、六官、隋的州县与租调。名词确实重要，因为它们让我们知道一个朝廷希望怎样分类人、地、役和官；但名词不是一台已经铺开的机器。诏书从都城出发，要经过中书、州郡、县吏、里正、豪强、寺院、军镇和家户，才可能在一块田、一项差役或一桩争讼中显形。每一层都有自己的信息、关系和资源，因此同一制度在不同地点会留下不同形状。

理解这一点并不等于把所有制度都说成“毫无作用”。恰恰相反，制度之所以值得追踪，是因为它改变了人们可以提出什么要求、官员必须保存何种文书、地方中介可以借什么名义行事。西晋编成律令，北魏颁布均田与三长，北周重整军府，隋整合州郡，都是为了把不稳定的征发、裁判和供给变成可重复的程序。问题在于：这些程序何时能把成本从临时暴力转移为常规负担，又何时反而暴露了国家看不见、管不到或无力支付的部分。

\subsection{西晋的统一：接收比宣布更难}

太康元年孙吴降服后，西晋拥有自三国以来少见的全国性名义。王濬舰队由长江而下，建业开城，使朝廷可以把吴的州郡、官员和百姓纳入晋的年号与官制。然而，军事胜利只解决了最高政权的归属。吴地已有自己的地方行政、军队遗存、豪强关系和水路经济；接收者必须决定哪些官继续留用、哪些仓储和船只可供调用、哪些旧有赋役要改写。若只把统一理解为孙皓出降的一刻，便看不见后来任何财政、地方治理或江南社会问题的起点。

西晋的律令与官制向天下提出一套可预期的裁判和任用语言。它们让中央能够说某件事合法、某种官职应由何等人担任，也为地方吏员提供了处理文书的格式。但法典文本本身不能告诉我们诉讼在每一县如何结束，或权贵、宗室和军队何时能绕过程序。八王之乱尤其说明，当持有兵权的宗王把军队带入洛阳时，顾命、诏令与官署并不足以保护政治秩序。制度不是被“废除”后才失灵；它可以继续被引用，却在武装力量、宫廷信息和粮饷都被私人集团控制时失去约束力。

地方社会也不只是等待中央命令的空白。坞壁、豪强、宗族和地方吏在危机中能组织自保、收留或排斥流民、保存粮食与道路消息。对居民而言，这类中介有时是避开战乱的条件，有时又意味着额外依附与征发。朝廷在安定时期需要借重它们以取得税粮和治安，战乱时则可能发现它们已拥有独立的武装与谈判能力。西晋末年的崩解并非由一个抽象“门阀”自动造成，而是中央财政、宗室兵权、边地压力与地方资源在特定时刻不再能被同一套程序接合。

\subsection{从品评到登记：国家怎样寻找可以承担责任的人}

九品中正制常被当成士族政治的速写。它确实让门第、乡里声望与官员选任发生稳定联系，但它不能把全部官僚与全部社会冻结为若干永恒家族。中正对人物的品评需要地方消息和既有关系；朝廷对这种消息的依赖既能减少陌生人的不确定性，也会使资源集中于能被识别、能相互担保的群体。南渡以后，旧郡望在江南仍有价值，新的仕宦与婚姻却必须适应当地社会。制度在这里保存了来处，也制造了进入官场的门槛。

户籍是另一种更贴近日常的寻找方式。国家并不直接面对“人口”，而是面对写在簿册上的户、丁、田、役与保。登记需要有人报告出生、死亡、迁居和土地占有，需要有人能到村落核验，也需要朝廷愿意承认某些例外。侨置与土断的反复正说明这个过程没有天然终点：寄居者可能保有旧籍，地方可能不愿失去劳力，豪强可能把依附者留在看不清的关系中，战乱又会使前一年登记立刻过时。户籍的价值在于让征收和赈济有了可操作对象；它的危险在于把复杂的家庭、迁徙和依附关系压进少数类别，并让未被正确登记的人更难主张权利。

北魏的制度转型应放在这一长期问题中看。平城朝廷在统一北方的过程中继承并改造郡县、军镇和部众关系；它需要把战争所获得的土地与人力转化为持续的供给。太武帝时期的征服可以扩大直接干预的区域，却也带来更多需要安置、分类和征调的人。把北魏的改革简单称作“汉化”会失去问题的重心：朝廷并不是在两种文化之间做一道选择题，而是在广阔、战后且社会关系多样的地区，寻找让税、役、兵和治安能够延续的中介。

\subsection{均田与三长：条文如何碰到田界与里门}

太和九年颁布均田令、次年推行三长制，是北魏行政史的关键节点。制度把土地的类别、受田者的责任、邻里编组与地方治安放进同一套设计：国家希望依据身份和劳力授受土地，以三长掌握户口、催课赋役，并减少强宗与旧部众对基层的独占。它是一项雄心极大的尝试，因为它承认没有可核验的家户和地方负责人，中央就无法把北方统一转化为稳定财政。

但必须区分诏令所说的规则与每一地真正发生的分配。田土有肥瘠和水旱之别，战争与迁徙使占有关系经常中断，地方强者和寺院可能已有难以重新划分的资源；三长若要履职，也必须依赖乡里认可、文书能力和上级支持。出土文书、墓志或个别地区的材料能帮助我们看到制度接触地面的片段，却不能证明全境同步丈量或完全均分。最稳妥的结论是：北魏确实以这些名目重组了国家与基层相互看见的方式，执行的密度和效果则因地区、时段和中介而异。

\begin{turningpoint}[均田与三长：把土地、赋役和治安绑成同一条责任链]
改革的创新不只是“给田”。它把谁可受田、谁应纳租调、邻里由谁申报和地方治安由谁负责放进相互牵连的制度。这样做能使朝廷减少临时搜括，代价是基层负责人获得新的权力，家户的迁居与隐匿也会成为围绕登记展开的博弈。改革能否成功，取决于田界、簿册和地方关系能否同时被维持。
\end{turningpoint}

班禄与官僚薪给的问题也不应从均田中抽离。若地方官与基层人员缺少稳定报酬，他们便更可能把征收、罚款或人情转化为个人资源；若国家能提供俸给，就有机会把部分取用纳入可监督的制度。可“有俸”不等于贪腐消失，“有法”也不等于地方弱者就能获得同等裁判。制度得失应问得更具体：它减少了哪一种任意征敛，把哪些记录和劳作加在谁身上，又为后来的地方中介留下什么新的机会。

\subsection{军府与州郡：兵从哪里来，粮又从哪里来}

六镇起事和北魏分裂之后，兵役与土地的关系更难维持在一张表格里。东魏、北齐的河北网络拥有城市、军队和旧镇人群，西魏、北周则在关中以较少的人口和不同的军政编组求存。后世常以“府兵制”概括西魏北周的军队组织，但若把它写成人人平等的农兵制度，便会忽略军户、将领、地方豪强和国家供给之间的等级与变化。军人能否在和平时生产、战时出征，取决于土地、家属、训练、武器和统帅网络是否都能被维系。

关中政权的选择尤其带有地理约束。它没有河北那样广大的资源腹地，却有山河形势、旧都遗产和可集中调度的核心区域。宇文泰及其继承者以军府、六官等安排整合军政，并非照搬古制的仪式；古名提供了正统语言，真正的问题仍是如何让关中有限的人力、粮食和官僚服务于与东面的长期竞争。北齐方面依靠河北的城市、手工业和人口，同样不能只被写成“富而不治”；强大的资源也会吸引宫廷、军人和地方集团争夺分配权。

北周灭北齐以后，接收邺及河北不只是增添地图面积。新政权需要处理旧军队、官员、户籍、仓储和宗教资源，既要利用已有生产与道路，又担心旧网络成为反抗的依托。隋取代北周、随后南下灭陈时也面对同样的接收技术。再统一的行政成就，不在于一个皇帝登基时说自己拥有天下，而在于州郡、县、仓、驿和簿册能否让不同区域的责任变得可追问。五八九年以后，江南没有立即变成关中；隋朝只能在接收中不断调整。

\subsection{法律、赦令与地方裁判}

法律史最容易把读者带进另一个误解：仿佛只要法典存在，社会就已经依照法典运转。两晋南北朝的律令、格式、赦令和官府文书当然为裁判提供了语言，也限制了某些官员可以公开宣称的理由；但案件是否能送到官府、证人能否作证、财物能否追回、强者是否可以施压，仍由地方条件决定。战乱和迁徙使证据、人证与原籍常常失散，恰使地方中介的作用更大。

赦令尤其能显出法与政治的交错。新君即位、迁都、战胜或灾异后，朝廷可以通过赦令重申宽宥与权威；对许多被牵连者，它可能带来真实的命运转折。但赦令并不是所有案件和所有区域的自动橡皮擦，执行者要解释条款，受害者仍要面对失去的土地、亲属与信誉。史书会记录宏大的赦天下，较少记录一县一里的具体余波。写制度时保留这种不对称，不是贬低法律，而是承认法律必须借由看得见的人和文书才会发生作用。

\begin{sourcenote}[制度材料能说到哪里]
正史的志、诏令与后出的编年材料最适合确认制度何时被宣布、朝廷怎样为其辩护；它们不能自动证明各地执行。墓志、契约、户籍残卷、造像题记和遗址可把某些家户、地方官或资源放在具体地点，却同样不能代表全境。把这些材料放在一起，能够描出制度从中央语言走向地方实践的路径；无法接上的部分，应保留为待问的问题。
\end{sourcenote}

从西晋接收吴地到隋接收陈地，制度不断重申同一种抱负：让统治不必每次都靠临时征兵、赏赐或掠夺，而能透过登记、任命、征收与裁判延续。它每次都要付出成本。更多登记意味着更多人要被看见和承担，更多中介意味着更多人可以操纵信息，更多统一标准又可能与地方习惯、地理和既有关系冲突。正因如此，制度不是王朝兴衰的背景板，而是人们在更替之间争取安全、资源和尊严的现场。
